% GENERATED FILE. Edit canonical lessons, then run npm run generate:books.
% canonical-source-hash: fnv1a64:2d36a5e4929d783d
% canonical-lessons: TE-C397-antekakunda, TE-C397-adananga, TE-C397-sanksiptanga, TE-C397-perkonu, TE-C397-vyaktam-ceyu, TE-R397-review-390-397

\chapter{Adding a Point}
\label{ch:te-adding-a-point}
\hlchaptermodality{397}

\begin{chapteropening}
\textbf{By the end of this chapter:} \emph{I can add points, sum up and say what I think.}

Complete the last lesson of chapter 397: I can add points, sum up and say what I think.
\end{chapteropening}

\section[\texorpdfstring{\te{అంతేకాకుండా} (antēkākuṇḍā)}{అంతేకాకుండా (antēkākuṇḍā)}]{\te{అంతేకాకుండా} (antēkākuṇḍā) — not only that; moreover}

\label{lesson:TE-C397-antekakunda}



\begin{quote}
Before the new one: say the Telugu for contrary, then the Telugu for likewise.
\end{quote}

\subsection*{You'll want to know: \te{అంతేకాకుండా}}
\textbf{\te{అంతేకాకుండా}} — \emph{antēkākuṇḍā} — \textquotedblleft{}not only that; moreover\textquotedblright{}.

\begin{quote}
\te{ఆ} \te{ఇల్లు} \te{పెద్దది}, \te{అంతేకాకుండా} \te{చౌక} \te{కూడా}.
\end{quote}

\begin{quote}
\emph{ā illu peddadi, antēkākuṇḍā cauka kūḍā.}
\end{quote}

\begin{quote}
That house is big; not only that, it is cheap too.
\end{quote}

\te{అంతే}, only that, + \te{కాకుండా}, without being: not only that.

\begin{cousinweb}
\textbf{\te{అంతేకాకుండా}} (\emph{antēkākuṇḍā}) is a native Telugu phrase: \emph{antē}, \textquotedblleft{}only that\textquotedblright{}, on the far-pointing base \emph{a-} (DEDR 1), and \emph{kākuṇḍā}, \textquotedblleft{}without being\textquotedblright{}, a form of \emph{avu}, \textquotedblleft{}to be\textquotedblright{} (DEDR 333).
\end{cousinweb}

\begin{grammarlens}[title={What you've built}]
One more word for giving reasons and opinions.
\end{grammarlens}

\subsection*{Your turn}
\begin{itemize}
\raggedright
  \item \emph{Say it:} \emph{antēkākuṇḍā}
  \item \emph{Say it:} \emph{antēkākuṇḍā}, once more
  \item \emph{Say it:} \emph{ā illu peddadi, antēkākuṇḍā cauka kūḍā.} — That house is big; not only that, it is cheap too.
  \item \emph{Recall:} say the Telugu for contrary, then the Telugu for likewise, then say \emph{antēkākuṇḍā} again
\end{itemize}

\begin{tcolorbox}[breakable,colback=teal!4,colframe=teal!35!black,title={Before you move on}]
What does \te{అంతేకాకుండా} mean? (\textquotedblleft{}Not only that; moreover\textquotedblright{}.) Say it once more.
\end{tcolorbox}
\section[\texorpdfstring{\te{అదనంగా} (adanaṅgā)}{అదనంగా (adanaṅgā)}]{\te{అదనంగా} (adanaṅgā) — in addition; extra}

\label{lesson:TE-C397-adananga}



\begin{quote}
Before the new one: say the Telugu for likewise, then the Telugu for not only that.
\end{quote}

\subsection*{You'll want to know: \te{అదనంగా}}
\textbf{\te{అదనంగా}} — \emph{adanaṅgā} — \textquotedblleft{}in addition; extra\textquotedblright{}.

\begin{quote}
\te{ఆదివారం} \te{పని} \te{చేస్తే} \te{అదనంగా} \te{డబ్బు} \te{ఇస్తారు}.
\end{quote}

\begin{quote}
\emph{ādivāraṁ pani cēstē adanaṅgā ḍabbu istāru.}
\end{quote}

\begin{quote}
If you work on Sunday they pay extra.
\end{quote}

Before a noun: \te{అదనపు} \te{ఖర్చు}, an extra cost.

\begin{cousinweb}
\textbf{\te{అదనంగా}} (\emph{adanaṅgā}) is native Dravidian: DEDR 136 lists Telugu \emph{adanamu}, \textquotedblleft{}exceeding, much, excessive\textquotedblright{}; -\te{గా} makes the adverb.
\end{cousinweb}

\begin{grammarlens}[title={What you've built}]
One more word for giving reasons and opinions.
\end{grammarlens}

\subsection*{Your turn}
\begin{itemize}
\raggedright
  \item \emph{Say it:} \emph{adanaṅgā}
  \item \emph{Say it:} \emph{adanaṅgā}, once more
  \item \emph{Say it:} \emph{ādivāraṁ pani cēstē adanaṅgā ḍabbu istāru.} — If you work on Sunday they pay extra.
  \item \emph{Recall:} say the Telugu for likewise, then the Telugu for not only that, then say \emph{adanaṅgā} again
\end{itemize}

\begin{tcolorbox}[breakable,colback=teal!4,colframe=teal!35!black,title={Before you move on}]
What does \te{అదనంగా} mean? (\textquotedblleft{}In addition; extra\textquotedblright{}.) Say it once more.
\end{tcolorbox}
\section[\texorpdfstring{\te{సంక్షిప్తంగా} (saṅkṣiptaṅgā)}{సంక్షిప్తంగా (saṅkṣiptaṅgā)}]{\te{సంక్షిప్తంగా} (saṅkṣiptaṅgā) — in short; briefly}

\label{lesson:TE-C397-sanksiptanga}



\begin{quote}
Before the new one: say the Telugu for not only that, then the Telugu for in addition.
\end{quote}

\subsection*{You'll want to know: \te{సంక్షిప్తంగా}}
\textbf{\te{సంక్షిప్తంగా}} — \emph{saṅkṣiptaṅgā} — \textquotedblleft{}in short; briefly\textquotedblright{}.

\begin{quote}
\te{సంక్షిప్తంగా} \te{చెప్పాలంటే}, \te{మా} \te{దగ్గర} \te{డబ్బు} \te{లేదు}.
\end{quote}

\begin{quote}
\emph{saṅkṣiptaṅgā ceppālaṇṭē, mā daggara ḍabbu lēdu.}
\end{quote}

\begin{quote}
To put it briefly, we have no money.
\end{quote}

\te{సంక్షిప్తంగా} \te{చెప్పాలంటే}, to put it in a few words.

\begin{cousinweb}
\textbf{\te{సంక్షిప్తంగా}} (\emph{saṅkṣiptaṅgā}) is built on Sanskrit \emph{saṁkṣipta}, \textquotedblleft{}abbreviated, contracted, condensed\textquotedblright{} (Monier-Williams), with the Telugu adverb ending -\te{గా}.
\end{cousinweb}

\begin{grammarlens}[title={What you've built}]
One more word for giving reasons and opinions.
\end{grammarlens}

\subsection*{Your turn}
\begin{itemize}
\raggedright
  \item \emph{Say it:} \emph{saṅkṣiptaṅgā}
  \item \emph{Say it:} \emph{saṅkṣiptaṅgā}, once more
  \item \emph{Say it:} \emph{saṅkṣiptaṅgā ceppālaṇṭē, mā daggara ḍabbu lēdu.} — To put it briefly, we have no money.
  \item \emph{Recall:} say the Telugu for not only that, then the Telugu for in addition, then say \emph{saṅkṣiptaṅgā} again
\end{itemize}

\begin{tcolorbox}[breakable,colback=teal!4,colframe=teal!35!black,title={Before you move on}]
What does \te{సంక్షిప్తంగా} mean? (\textquotedblleft{}In short; briefly\textquotedblright{}.) Say it once more.
\end{tcolorbox}
\section[\texorpdfstring{\te{పేర్కొను} (pērkonu)}{పేర్కొను (pērkonu)}]{\te{పేర్కొను} (pērkonu) — to mention; to state}

\label{lesson:TE-C397-perkonu}



\begin{quote}
Before the new one: say the Telugu for in addition, then the Telugu for in short.
\end{quote}

\subsection*{You'll want to know: \te{పేర్కొను}}
\textbf{\te{పేర్కొను}} — \emph{pērkonu} — \textquotedblleft{}to mention; to state\textquotedblright{}.

\begin{quote}
\te{ఆ} \te{ఉత్తరంలో} \te{అతను} \te{తన} \te{కారణాలు} \te{పేర్కొన్నాడు}.
\end{quote}

\begin{quote}
\emph{ā uttaraṁlō atanu tana kāraṇālu pērkonnāḍu.}
\end{quote}

\begin{quote}
In that letter he stated his reasons.
\end{quote}

Past (I …): \te{పేర్కొన్నాను}, \emph{pērkonnānu}.

Formal and written; in talk, \te{చెప్పు}.

\begin{cousinweb}
\textbf{\te{పేర్కొను}} (\emph{pērkonu}) is native Dravidian, a compound of \emph{pēru}, \textquotedblleft{}a name\textquotedblright{} (DEDR 4410), and \emph{konu}, \textquotedblleft{}to take\textquotedblright{} (DEDR 2151): to name a thing.
\end{cousinweb}

\begin{grammarlens}[title={What you've built}]
One more word for giving reasons and opinions.
\end{grammarlens}

\subsection*{Your turn}
\begin{itemize}
\raggedright
  \item \emph{Say it:} \emph{pērkonu}
  \item \emph{Say it:} \emph{pērkonu}, once more
  \item \emph{Say it:} \emph{ā uttaraṁlō atanu tana kāraṇālu pērkonnāḍu.} — In that letter he stated his reasons.
  \item \emph{Recall:} say the Telugu for in addition, then the Telugu for in short, then say \emph{pērkonu} again
\end{itemize}

\begin{tcolorbox}[breakable,colback=teal!4,colframe=teal!35!black,title={Before you move on}]
What does \te{పేర్కొను} mean? (\textquotedblleft{}To mention; to state\textquotedblright{}.) Say it once more.
\end{tcolorbox}
\section[\texorpdfstring{\te{వ్యక్తం} \te{చేయు} (vyaktaṁ cēyu)}{వ్యక్తం చేయు (vyaktaṁ cēyu)}]{\te{వ్యక్తం} \te{చేయు} (vyaktaṁ cēyu) — to express; to voice}

\label{lesson:TE-C397-vyaktam-ceyu}



\begin{quote}
Before the new one: say the Telugu for in short, then the Telugu for to mention.
\end{quote}

\subsection*{You'll want to know: \te{వ్యక్తం} \te{చేయు}}
\textbf{\te{వ్యక్తం} \te{చేయు}} — \emph{vyaktaṁ cēyu} — \textquotedblleft{}to express; to voice\textquotedblright{}.

\begin{quote}
\te{ఆమె} \te{తన} \te{సంతోషాన్ని} \te{వ్యక్తం} \te{చేసింది}.
\end{quote}

\begin{quote}
\emph{āme tana santōṣānni vyaktaṁ cēsindi.}
\end{quote}

\begin{quote}
She expressed her happiness.
\end{quote}

Past (I …): \te{వ్యక్తం} \te{చేశాను}, \emph{vyaktaṁ cēśānu}.

For feelings and views: \te{అభిప్రాయం} \te{వ్యక్తం} \te{చేయు}, to express an opinion.

\begin{cousinweb}
\textbf{\te{వ్యక్తం} \te{చేయు}} (\emph{vyaktaṁ cēyu}) is built from Sanskrit \emph{vyakta}, \textquotedblleft{}manifested, apparent, evident\textquotedblright{} (Monier-Williams), and the native verb \te{చేయు}, \textquotedblleft{}to do\textquotedblright{} (DEDR 1957): to make a feeling plain.
\end{cousinweb}

\begin{grammarlens}[title={What you've built}]
One more word for giving reasons and opinions.
\end{grammarlens}

\subsection*{Your turn}
\begin{itemize}
\raggedright
  \item \emph{Say it:} \emph{vyaktaṁ cēyu}
  \item \emph{Say it:} \emph{vyaktaṁ cēyu}, once more
  \item \emph{Say it:} \emph{āme tana santōṣānni vyaktaṁ cēsindi.} — She expressed her happiness.
  \item \emph{Recall:} say the Telugu for in short, then the Telugu for to mention, then say \emph{vyaktaṁ cēyu} again
\end{itemize}

\begin{tcolorbox}[breakable,colback=teal!4,colframe=teal!35!black,title={Before you move on}]
What does \te{వ్యక్తం} \te{చేయు} mean? (\textquotedblleft{}To express; to voice\textquotedblright{}.) Say it once more.
\end{tcolorbox}
\section[(dialogue)]{Words from chapters 390–397 — a review}

\label{lesson:TE-R397-review-390-397}



\begin{quote}
Say the Telugu for a fact and a theory.
\end{quote}

\subsection*{Your turn}
Say these aloud:

\begin{itemize}
\raggedright
  \item a fact, a theory, a principle, a process, to define
  \item importance, priority, basic, to stress, to bring up
  \item a source, to originate, to lead to, to arise, to crop up
  \item a benefit, to benefit, to suffer a loss, harmful, to be useful
  \item an explanation, clarity, a justification, to justify, exact
  \item so, for that reason, in view of, that is to say, a footing
  \item nevertheless, on the other hand, whereas, contrary, likewise
  \item not only that, in addition, in short, to mention, to express
\end{itemize}

\begin{tcolorbox}[breakable,colback=teal!4,colframe=teal!35!black,title={Before you move on}]
A fact? (\textbf{\te{వాస్తవం}}, \emph{vāstavaṁ}.) A theory? (\textbf{\te{సిద్ధాంతం}}, \emph{siddhāntaṁ}.) A principle? (\textbf{\te{సూత్రం}}, \emph{sūtraṁ}.) A process? (\textbf{\te{ప్రక్రియ}}, \emph{prakriya}.) To define? (\textbf{\te{నిర్వచించు}}, \emph{nirvaciñcu}.) Importance? (\textbf{\te{ప్రాముఖ్యత}}, \emph{prāmukhyata}.) Priority? (\textbf{\te{ప్రాధాన్యం}}, \emph{prādhānyaṁ}.) Basic? (\textbf{\te{ప్రాథమిక}}, \emph{prāthamika}.) To stress? (\textbf{\te{నొక్కిచెప్పు}}, \emph{nokkiceppu}.) To bring up? (\textbf{\te{ప్రస్తావించు}}, \emph{prastāviñcu}.) A source? (\textbf{\te{మూలం}}, \emph{mūlaṁ}.) To originate? (\textbf{\te{ఉద్భవించు}}, \emph{udbhaviñcu}.) To lead to? (\textbf{\te{దారితీయు}}, \emph{dāritīyu}.) To arise? (\textbf{\te{కలుగు}}, \emph{kalugu}.) To crop up? (\textbf{\te{తలెత్తు}}, \emph{talettu}.) A benefit? (\textbf{\te{ప్రయోజనం}}, \emph{prayōjanaṁ}.) To benefit? (\textbf{\te{లాభపడు}}, \emph{lābhapaḍu}.) To suffer a loss? (\textbf{\te{నష్టపోవు}}, \emph{naṣṭapōvu}.) Harmful? (\textbf{\te{హానికరమైన}}, \emph{hānikaramaina}.) To be useful? (\textbf{\te{ఉపయోగపడు}}, \emph{upayōgapaḍu}.) An explanation? (\textbf{\te{వివరణ}}, \emph{vivaraṇa}.) Clarity? (\textbf{\te{స్పష్టత}}, \emph{spaṣṭata}.) A justification? (\textbf{\te{సమర్థన}}, \emph{samarthana}.) To justify? (\textbf{\te{సమర్థించు}}, \emph{samarthiñcu}.) Exact? (\textbf{\te{ఖచ్చితమైన}}, \emph{khaccitamaina}.) So? (\textbf{\te{కనుక}}, \emph{kanuka}.) For that reason? (\textbf{\te{అందుచేత}}, \emph{anducēta}.) In view of? (\textbf{\te{దృష్ట్యా}}, \emph{dṛṣṭyā}.) That is to say? (\textbf{\te{అనగా}}, \emph{anagā}.) A footing? (\textbf{\te{ప్రాతిపదిక}}, \emph{prātipadika}.) Nevertheless? (\textbf{\te{అయినప్పటికీ}}, \emph{ayinappaṭikī}.) On the other hand? (\textbf{\te{మరోవైపు}}, \emph{marōvaipu}.) Whereas? (\textbf{\te{కాగా}}, \emph{kāgā}.) Contrary? (\textbf{\te{విరుద్ధం}}, \emph{viruddhaṁ}.) Likewise? (\textbf{\te{అదేవిధంగా}}, \emph{adēvidhaṅgā}.) Not only that? (\textbf{\te{అంతేకాకుండా}}, \emph{antēkākuṇḍā}.) In addition? (\textbf{\te{అదనంగా}}, \emph{adanaṅgā}.) In short? (\textbf{\te{సంక్షిప్తంగా}}, \emph{saṅkṣiptaṅgā}.) To mention? (\textbf{\te{పేర్కొను}}, \emph{pērkonu}.) To express? (\textbf{\te{వ్యక్తం} \te{చేయు}}, \emph{vyaktaṁ cēyu}.)
\end{tcolorbox}
